% ============================================================
% Question Bank: ch02-08 Personal Financial Literacy
% Topic Title: Personal Financial Literacy
% CCSS: Supplementary
% Questions: 27 (15 MC + 6 SA + 6 Graphical)
% ============================================================

% --- Multiple Choice Questions (15) ---

\begin{practiceQuestion}{2-8-q01}{mc}
\begin{questionText}
Which type of financial institution is a member-owned cooperative that typically offers lower fees?
\end{questionText}
\choiceA{Investment firm}
\choiceB{Credit union}
\choiceC{Insurance company}
\choiceD{Stock exchange}
\correctAnswer{B}
\explanation{Credit unions are member-owned cooperatives that often provide lower fees and better interest rates compared to traditional banks.}
\end{practiceQuestion}

\begin{practiceQuestion}{2-8-q02}{mc}
\begin{questionText}
A debit card purchase of $\$45$ is made from an account with a $\$200$ balance. What is the new balance?
\end{questionText}
\choiceA{$\$155$}
\choiceB{$\$200$}
\choiceC{$\$245$}
\choiceD{$\$45$}
\correctAnswer{A}
\explanation{A debit card takes money directly from your account: $200 - 45 = \$155$.}
\end{practiceQuestion}

\begin{practiceQuestion}{2-8-q03}{mc}
\begin{questionText}
A credit card has a $\$500$ balance and charges $1.5\%$ interest per month. How much interest is charged the first month?
\end{questionText}
\choiceA{$\$1.50$}
\choiceB{$\$5.00$}
\choiceC{$\$7.50$}
\choiceD{$\$75.00$}
\correctAnswer{C}
\explanation{Monthly interest: $500 \times 0.015 = \$7.50$.}
\end{practiceQuestion}

\begin{practiceQuestion}{2-8-q04}{mc}
\begin{questionText}
Which of the following is a disadvantage of using a credit card instead of a debit card?
\end{questionText}
\choiceA{It does not require a bank account}
\choiceB{You may be charged interest on unpaid balances}
\choiceC{Purchases are limited to your checking balance}
\choiceD{There is no purchase protection}
\correctAnswer{B}
\explanation{Credit cards charge interest on any unpaid balance, which can increase the total cost. Debit cards draw directly from your account with no interest.}
\end{practiceQuestion}

\begin{practiceQuestion}{2-8-q05}{mc}
\begin{questionText}
Bank A charges a $\$3$ monthly maintenance fee. Bank B charges no fee but requires a $\$500$ minimum balance. If you keep $\$350$ in your account all year, how much would Bank A cost you in fees per year?
\end{questionText}
\choiceA{$\$3$}
\choiceB{$\$12$}
\choiceC{$\$30$}
\choiceD{$\$36$}
\correctAnswer{D}
\explanation{Monthly fee: $\$3 \times 12 = \$36$ per year.}
\end{practiceQuestion}

\begin{practiceQuestion}{2-8-q06}{mc}
\begin{questionText}
A checking account earns $0.5\%$ annual interest on a $\$2{,}000$ balance. How much interest is earned in one year?
\end{questionText}
\choiceA{$\$1$}
\choiceB{$\$5$}
\choiceC{$\$10$}
\choiceD{$\$100$}
\correctAnswer{C}
\explanation{$2{,}000 \times 0.005 = \$10$.}
\end{practiceQuestion}

\begin{practiceQuestion}{2-8-q07}{mc}
\begin{questionText}
If you do not pay your full credit card balance each month, which of the following happens?
\end{questionText}
\choiceA{The bank closes your account}
\choiceB{You earn rewards faster}
\choiceC{Interest is charged on the remaining balance}
\choiceD{Your credit limit increases automatically}
\correctAnswer{C}
\explanation{Credit card companies charge interest on any balance not paid by the due date each month.}
\end{practiceQuestion}

\begin{practiceQuestion}{2-8-q08}{mc}
\begin{questionText}
A credit card balance of $\$800$ charges $18\%$ annual interest. Approximately how much interest is charged per month?
\end{questionText}
\choiceA{$\$8$}
\choiceB{$\$12$}
\choiceC{$\$14.40$}
\choiceD{$\$144$}
\correctAnswer{B}
\explanation{Monthly rate: $18\% \div 12 = 1.5\%$. Monthly interest: $800 \times 0.015 = \$12$.}
\end{practiceQuestion}

\begin{practiceQuestion}{2-8-q09}{mc}
\begin{questionText}
Which payment method helps you avoid spending more money than you have?
\end{questionText}
\choiceA{Credit card}
\choiceB{Store credit line}
\choiceC{Debit card}
\choiceD{Personal loan}
\correctAnswer{C}
\explanation{A debit card is linked directly to your bank account, so you can only spend what you have.}
\end{practiceQuestion}

\begin{practiceQuestion}{2-8-q10}{mc}
\begin{questionText}
A savings account at a bank pays $2\%$ annual interest. A certificate of deposit (CD) at the same bank pays $4\%$ annual interest but locks the money for $1$ year. On a $\$1{,}000$ deposit, how much more would the CD earn in one year?
\end{questionText}
\choiceA{$\$2$}
\choiceB{$\$10$}
\choiceC{$\$20$}
\choiceD{$\$40$}
\correctAnswer{C}
\explanation{Savings: $1{,}000 \times 0.02 = \$20$. CD: $1{,}000 \times 0.04 = \$40$. Difference: $40 - 20 = \$20$.}
\end{practiceQuestion}

\begin{practiceQuestion}{2-8-q11}{mc}
\begin{questionText}
An ATM charges a $\$2.50$ fee each time you withdraw from another bank's machine. If you make $8$ withdrawals per month from another bank's ATM, how much do you pay in fees per year?
\end{questionText}
\choiceA{$\$20$}
\choiceB{$\$30$}
\choiceC{$\$120$}
\choiceD{$\$240$}
\correctAnswer{D}
\explanation{Monthly fees: $2.50 \times 8 = \$20$. Annual: $20 \times 12 = \$240$.}
\end{practiceQuestion}

\begin{practiceQuestion}{2-8-q12}{mc}
\begin{questionText}
Which of the following is NOT typically a service offered by a bank?
\end{questionText}
\choiceA{Checking accounts}
\choiceB{Savings accounts}
\choiceC{Home insurance policies}
\choiceD{Personal loans}
\correctAnswer{C}
\explanation{Banks offer checking and savings accounts and loans, but home insurance is typically sold by insurance companies.}
\end{practiceQuestion}

\begin{practiceQuestion}{2-8-q13}{mc}
\begin{questionText}
A credit card has a $\$1{,}200$ balance. You pay $\$100$ each month and the card charges $1\%$ monthly interest on the remaining balance. After your first payment, what is the new balance before the next payment?
\end{questionText}
\choiceA{$\$1{,}100$}
\choiceB{$\$1{,}111$}
\choiceC{$\$1{,}112$}
\choiceD{$\$1{,}200$}
\correctAnswer{B}
\explanation{After payment: $1{,}200 - 100 = \$1{,}100$. Interest: $1{,}100 \times 0.01 = \$11$. New balance: $1{,}100 + 11 = \$1{,}111$.}
\end{practiceQuestion}

\begin{practiceQuestion}{2-8-q14}{mc}
\begin{questionText}
What is one advantage of using a credit card responsibly?
\end{questionText}
\choiceA{You never have to pay back the money}
\choiceB{It builds a positive credit history}
\choiceC{It increases your bank savings}
\choiceD{Interest is never charged}
\correctAnswer{B}
\explanation{Responsible use of a credit card (paying on time and in full) helps build a positive credit history, which impacts borrowing in the future.}
\end{practiceQuestion}

\begin{practiceQuestion}{2-8-q15}{mc}
\begin{questionText}
A bank account charges a $\$35$ overdraft fee each time the balance goes below $\$0$. If someone overdrafts $3$ times in one month, how much is charged in fees?
\end{questionText}
\choiceA{$\$35$}
\choiceB{$\$70$}
\choiceC{$\$105$}
\choiceD{$\$350$}
\correctAnswer{C}
\explanation{$35 \times 3 = \$105$ in overdraft fees.}
\end{practiceQuestion}

% --- Short Answer Questions (6) ---

\begin{practiceQuestion}{2-8-q16}{sa}
\begin{questionText}
A credit card charges $2\%$ monthly interest. If you carry a $\$600$ balance, how much interest is charged in one month?
\end{questionText}
\correctAnswer{$\$12$}
\explanation{$600 \times 0.02 = \$12$.}
\end{practiceQuestion}

\begin{practiceQuestion}{2-8-q17}{sa}
\begin{questionText}
You have $\$320$ in your checking account and use your debit card to buy $5$ items that each cost $\$72$. Is there enough money? If not, how much are you short?
\end{questionText}
\correctAnswer{Short $\$40$}
\explanation{Total: $5 \times 72 = \$360$. Balance: $\$320$. Since $360 > 320$, you are short $360 - 320 = \$40$.}
\end{practiceQuestion}

\begin{practiceQuestion}{2-8-q18}{sa}
\begin{questionText}
A bank pays $1.5\%$ annual interest on savings. How much interest does a $\$4{,}000$ balance earn in $2$ years?
\end{questionText}
\correctAnswer{$\$120$}
\explanation{$I = 4{,}000 \times 0.015 \times 2 = \$120$.}
\end{practiceQuestion}

\begin{practiceQuestion}{2-8-q19}{sa}
\begin{questionText}
A store credit card charges $21\%$ annual interest. How much interest is charged on a $\$300$ balance after one year?
\end{questionText}
\correctAnswer{$\$63$}
\explanation{$300 \times 0.21 = \$63$.}
\end{practiceQuestion}

\begin{practiceQuestion}{2-8-q20}{sa}
\begin{questionText}
A bank charges a $\$5$ monthly service fee. A credit union charges $\$2$ per month. How much more do you pay at the bank over $3$ years?
\end{questionText}
\correctAnswer{$\$108$}
\explanation{Bank: $5 \times 36 = \$180$. Credit union: $2 \times 36 = \$72$. Difference: $180 - 72 = \$108$.}
\end{practiceQuestion}

\begin{practiceQuestion}{2-8-q21}{sa}
\begin{questionText}
If a credit card has a $\$2{,}000$ limit and you owe $\$1{,}350$, how much available credit do you have?
\end{questionText}
\correctAnswer{$\$650$}
\explanation{$2{,}000 - 1{,}350 = \$650$ available credit.}
\end{practiceQuestion}

% --- Graphical Questions (3 GMC + 3 GSA) ---

\begin{practiceQuestion}{2-8-q22}{gmc}
\begin{questionText}
The table shows monthly fees for three bank accounts. Which account has the lowest total annual cost?

\begin{center}
\begin{tabular}{|l|c|c|}
\hline
\textbf{Account} & \textbf{Monthly Fee} & \textbf{ATM Fee (per use)} \\
\hline
Basic & $\$5$ & $\$0$ \\
\hline
Standard & $\$0$ & $\$3$ \\
\hline
Premium & $\$8$ & $\$0$ \\
\hline
\end{tabular}
\end{center}

Assume you use an ATM $2$ times per month.
\end{questionText}
\choiceA{Basic ($\$60$)}
\choiceB{Standard ($\$72$)}
\choiceC{Premium ($\$96$)}
\choiceD{Basic and Standard are tied}
\correctAnswer{A}
\explanation{Basic: $5 \times 12 = \$60$. Standard: $3 \times 2 \times 12 = \$72$. Premium: $8 \times 12 = \$96$. Basic is cheapest.}
\end{practiceQuestion}

\begin{practiceQuestion}{2-8-q23}{gmc}
\begin{questionText}
The bar graph compares interest charged on a $\$1{,}000$ balance at two different annual rates over one year. What is the monthly interest for Card B?

\begin{center}
\begin{tikzpicture}[scale=0.5]
  \draw[thick,->] (0,0) -- (0,6) node[above]{\small Annual Interest (\$)};
  \draw[thick,->] (0,0) -- (6,0);
  \fill[funBlue!60] (0.75,0) rectangle (2.25,3);
  \fill[funRed!60] (3.25,0) rectangle (4.75,4.5);
  \node[below,font=\small] at (1.5,0) {Card A};
  \node[below,font=\small] at (4,0) {Card B};
  \foreach \y/\lab in {1.5/150,3/300,4.5/450} {
    \draw (-0.1,\y) -- (0.1,\y) node[left,font=\small,xshift=-2pt] {$\$\lab$};
  }
  \node[above,font=\small] at (1.5,3) {$15\%$};
  \node[above,font=\small] at (4,4.5) {$22.5\%$};
\end{tikzpicture}
\end{center}
\end{questionText}
\choiceA{$\$12.50$}
\choiceB{$\$18.75$}
\choiceC{$\$22.50$}
\choiceD{$\$37.50$}
\correctAnswer{B}
\explanation{Card B annual: $1{,}000 \times 0.225 = \$225$. Monthly: $225 \div 12 = \$18.75$.}
\end{practiceQuestion}

\begin{practiceQuestion}{2-8-q24}{gmc}
\begin{questionText}
The table compares two payment methods for a $\$400$ purchase. Which method costs less overall and by how much?

\begin{center}
\begin{tabular}{|l|c|c|c|}
\hline
\textbf{Method} & \textbf{Purchase} & \textbf{Interest} & \textbf{Total Paid} \\
\hline
Debit Card & $\$400$ & $\$0$ & $\$400$ \\
\hline
Credit Card (paid over 6 months) & $\$400$ & $\$36$ & $\$436$ \\
\hline
\end{tabular}
\end{center}
\end{questionText}
\choiceA{Credit card by $\$36$}
\choiceB{Debit card by $\$36$}
\choiceC{They cost the same}
\choiceD{Debit card by $\$400$}
\correctAnswer{B}
\explanation{Debit: $\$400$ total. Credit: $\$436$ total. Debit saves $436 - 400 = \$36$.}
\end{practiceQuestion}

\begin{practiceQuestion}{2-8-q25}{gsa}
\begin{questionText}
The table shows a checking account register. What is the final balance?

\begin{center}
\begin{tabular}{|l|c|c|c|}
\hline
\textbf{Transaction} & \textbf{Deposit} & \textbf{Withdrawal} & \textbf{Balance} \\
\hline
Starting & & & $\$500$ \\
\hline
Paycheck & $\$350$ & & \\
\hline
Rent & & $\$400$ & \\
\hline
Groceries & & $\$75$ & \\
\hline
Birthday gift & $\$50$ & & \\
\hline
\end{tabular}
\end{center}
\end{questionText}
\correctAnswer{$\$425$}
\explanation{$500 + 350 - 400 - 75 + 50 = \$425$.}
\end{practiceQuestion}

\begin{practiceQuestion}{2-8-q26}{gsa}
\begin{questionText}
The pie chart shows how a $\$24$ monthly bank fee is divided. How much goes to the ATM fee category?

\begin{center}
\begin{tikzpicture}[scale=1.55]
  \fill[funBlue!50] (0,0) -- (90:1) arc (90:234:1) -- cycle;
  \fill[funGreen!50] (0,0) -- (234:1) arc (234:306:1) -- cycle;
  \fill[funOrange!50] (0,0) -- (306:1) arc (306:450:1) -- cycle;
  \draw (0,0) circle (1);
  \draw (0,0) -- (90:1); \draw (0,0) -- (234:1); \draw (0,0) -- (306:1);
  \node[font=\small,align=center] at (162:0.55) {Maint.\\$40\%$};
  \node[font=\small,align=center] at (270:0.48) {ATM\\$20\%$};
  \node[font=\small,align=center] at (18:0.55) {Other\\$40\%$};
\end{tikzpicture}
\end{center}
\end{questionText}
\correctAnswer{$\$4.80$}
\explanation{ATM portion: $20\%$ of $\$24 = 0.20 \times 24 = \$4.80$.}
\end{practiceQuestion}

\begin{practiceQuestion}{2-8-q27}{gsa}
\begin{questionText}
The table compares interest earned on $\$3{,}000$ at two banks over $3$ years. How much more does Bank B pay in interest?

\begin{center}
\begin{tabular}{|l|c|c|}
\hline
 & \textbf{Rate} & \textbf{Time} \\
\hline
Bank A & $1\%$ & $3$ years \\
\hline
Bank B & $3\%$ & $3$ years \\
\hline
\end{tabular}
\end{center}
\end{questionText}
\correctAnswer{$\$180$}
\explanation{Bank A: $3{,}000 \times 0.01 \times 3 = \$90$. Bank B: $3{,}000 \times 0.03 \times 3 = \$270$. Difference: $270 - 90 = \$180$.}
\end{practiceQuestion}
